% Abstract page (page i): title block + abstract + keywords, as in the example theses.
% TODO (author): revise once the pending results (paired cross-module tests, privacy
% measurements, stale-head cost) are in; keep it under one page.
\abstractheader
\addcontentsline{toc}{chapter}{Abstract}
\thispagestyle{plain}

Federated recommendation trains a collaborative-filtering model across many devices without
collecting the users' interaction histories on a server. Prior work agrees that some part of the
model has to be personalised, but disagrees on \emph{where}: nothing (all parameters global), the
user vector, a private scoring head, or a blend between a private part and population knowledge.
This thesis compares the four placements under one cross-device protocol on MovieLens 1M with
\DsUsers{} users, one user per client, a participation fraction of $C=\ParticipationC$ and
\NRounds{} rounds, scoring every user with the same leave-one-out protocol and 99 sampled
negatives.

The main finding is that the placement of personalisation dominates every other design choice in
this setting. Keeping the user vector on the device raises \ndcg{} by
\GainPersonalizedVsBaselinePct\% over the all-global baseline in an otherwise identical
configuration (\NdcgBaseline{} to \NdcgPersonalized), and a private scoring head on a shared item
table (PFedRec with final calibration) reaches \NdcgPfedrec. Replacing the private head by the
population-mean head costs \AbsDeltaHeadGlobalHead{} \ndcg{} on the same model, and interpolating
between the two with a data-driven weight does not recover the loss, so the value lies in having a
private head, not in blending it. The proposed adaptive method, which blends the user vector with a
population prototype and adds a per-user neural scorer, is reported as a negative result
(\NdcgAdaptive): its blending weight pulls user taste toward the mean, its per-user network has
\ParamsPersonalMlp{} parameters for a median of \DsMedianRatings{} interactions, and it uploads
the user vector, which weakens the privacy boundary.

Two further contributions concern how such comparisons should be evaluated. First, the gap between
the PFedRec number reported here and the published one is measured rather than assumed: scoring
with a client-side post-tuned item table lowers \ndcg{} (\NdcgEOneClean) when the test item can be
drawn as a training negative, and raises it (\NdcgEOneLeaky) only when the test item is excluded
from the negatives as in the official code, a difference of \PtContrastEOne{} concentrated in
\PtContrastEOneNonZero{} users. Second, stateful clients require evaluation safeguards: a
cold-state evaluation error that reduced \ndcg{} \ColdEvalRatio-fold was found and fixed, and an
offline re-scoring tool reproduces the federated evaluation exactly before any further analysis.
\pending{one sentence on the privacy measurements (byte-level communication table and the
preference-reconstruction attack from the uploaded user vector) once they are run}.

\medskip
\noindent\textbf{Keywords:} federated learning; recommender systems; collaborative filtering;
personalisation; privacy; evaluation protocol; MovieLens.
